%! Author = lili
%! Date = 8/20/26

\newglossaryentry{dreihundertnmfibrille}{
    name={300 nm Fibrille},
    description={Höhere Verpackungsebene des Chromatins, bei der die 30-nm-Chromatinfaser durch Ausbildung von Schleifen weiter zu einer etwa 300 nm dicken Struktur kondensiert wird.},
    type=defi
}


\newglossaryentry{zweifaktor_kreuzung}{
    name={Zweifaktor-Kreuzung},
    description={Kreuzungsexperiment, bei dem gleichzeitig die Vererbung zweier Merkmale (zwei Genorte) verfolgt wird, dient z. B. der Prüfung auf unabhängige Vererbung nach Mendel.},
    type=defi
}


\newglossaryentry{zentriole}{
    name={Zentriole},
    description={Zentralkörperchen; in der Nähe des
    Zellkerns gelegenes, meist paarweise (als Diplosom) vorkommendes Organell. Organisationszentrum für die
    Spindel. Die beiden Z.en sind rechtwinklig zueinander angeordnet, entsprechen in ihrem Aufbau den
    Basalkörpern von Zilien und Flagellen und sind wechselseitig in diese umwandelbar. Sind sind nur elektronen-
    mikroskopisch sichtbar.},
    type=defi
}

% TODO Noch Zentrosom?

\newglossaryentry{zweichrchr}{
    name={Zwei-Chromatid-Chromosom},
    description={Repliziertes Chromosom, das aus zwei genetisch identischen Schwesterchromatiden besteht, die am
    Centromer durch Cohesin zusammengehalten werden; entsteht nach der S-Phase.},
    type=defi
}
\newglossaryentry{xci}{
    name={X-Chromosom Inaktivierung},
    description={Prozess bei weiblichen Säugetieren, bei dem eines der beiden X-Chromosomen pro Zelle zufällig und stabil inaktiviert wird, um die Gendosis im Vergleich zu männlichen Individuen auszugleichen.},
    type=defi
}

\newglossaryentry{wobble}{
    name={Wobble-Position},
    description={Dritte Codonposition (3'), an der nicht-Watson-Crick-Basenpaarungen toleriert werden.
    Ermöglicht einer tRNA die Erkennung mehrerer synonymer Codons (z.\,B. G\,:\,U-Paar).},
    type=defi
}

\newglossaryentry{vntr}{
    name={Variable number tandem repeats},
    description={=Minisatellit},
    type=defi
}

\newglossaryentry{vererbung}{
    name={Vererbung},
    description={Weitergabe von Merkmalen (z.B. Augenfarbe) an die
    nachfolgende Generation},
    type=defi
}

\newglossaryentry{variabilität}{
    name={Variabilität},
    description={  Unterschiede der Nachkommen gegenüber ihren Eltern,
    auch Unterschiede der Mitglieder einer Population
    untereinander},
    type=defi
}


\newglossaryentry{translation}{
    name={Translation},
    description={Proteinsynthese
    an einer RNA-Matrize mittels Ribosomen},
    type=defi
}

\newglossaryentry{tetrade}{
    name={Tetrade},
    description={Gruppe der vier haploiden Zellen (Sporen), die aus einer Meiose einer diploiden Zelle hervorgehen.},
    type=defi
}



\newglossaryentry{tetradenanalyse}{
    name={Tetradenanalyse},
    description={Genetische Methode, bei der die vier Produkte einer einzelnen Meiose gemeinsam analysiert werden, um Kopplung, Rekombination oder Genkonversion nachzuweisen, vor allem bei Pilzen wie Neurospora.},
    type=defi
}


\newglossaryentry{transelement}{
    name={trans-Element},
    description={Transkriptionsfaktoren (Proteine), die an die cis-Elemente binden},
    type=defi
}




\newglossaryentry{Transkriptom}{
    name={Transkriptom},
    description={Gesamtheit aller (auf RNA-Ebene) exprimierten Gene Gesamtheit aller (auf RNA-Ebene) exprimierten
    Gene (inkl. mRNAs, tRNAs, rRNAs, ncRNAs, miRNAa, ...)},
    type=defi
}

\newglossaryentry{synapsis}{
    name={Synapsis},
    description={Paarung homologer Chromosomen},
    type=defi
}

\newglossaryentry{spontane_mutation}{
    name={spontane Mutation},
    description={Mutation, die ohne äußeren mutagenen Einfluss entsteht, z. B. durch Replikationsfehler oder tautomere Basenverschiebung.},
    type=defi
}

\newglossaryentry{spindel}{
    name={Spindel},
    description={Spindelförmige Anordnung von Mikrotubuli,
    die bei der Kernteilung die Chromosomen zu den beiden Spindelpolen (Zentriolen) transportieren.},
    type=defi
}
\newglossaryentry{smc}{
    name={SMC-Proteine},
    first={Structural Maintenance of Chromosomes Proteine},
    description={Familie chromosomaler ATPasen (Structural Maintenance of Chromosomes), beteiligt an Chromosomenkondensation, Schwesterchromatidkohäsion und DNA-Reparatur, z. B. Condensin und Cohesin.},
    type=defi
}

\newglossaryentry{sslp}{
    name={simple-sequence length polymorphism (SSLP)},
    description={Kurzer, in der Länge variabler, sich wiederholender DNA-Abschnitt (z. B. Mikrosatellit), dessen Längenpolymorphismus als genetischer Marker genutzt wird.},
    type=defi
}


\newglossaryentry{snp}{
    name={Single Nucleotid Polymorphism},
    description={Unterschiede in einzelnen Nukleotiden},
    type=defi
}


\newglossaryentry{str}{
    name={Short Tandem Repeat},
    description={=Mikrosatellit, SSR},
    type=defi
}

\newglossaryentry{sexFort}{
    name={sexuelle Fortpflanzung (generative, geschlechtliche)},
    description={ Erzeugung von Nachkommen, die eine Neukombination der
    Gene/Allele der Eltern tragen.
    Es entstehen individuelle und einzigartige Nachkommen.},
    type=defi
}

\newglossaryentry{schwesterchromosom}{
    name={Schwesterchromosom},
    description={todo},
    type=defi
}

\newglossaryentry{schwesterchromatide}{
    name={Schwesterchromatide},
    description={Durch identi-
    sche Verdopplung (Replikation) eines Chromosoms
    entstandene Chromatiden},
    type=defi
}

\newglossaryentry{rückkreuzung}{
    name={Rückkreuzung},
    description={
        Von einer Rückkreuzung spricht man, wenn schon einer der beiden Eltern für beide rezessiven Allele homozygot war;
        allgemein nennt man diesen Kreuzungstyp Testkreuzung.
    },
    type=defi
}

\newglossaryentry{ribonukleoprotein}{
    name={Ribonukleoprotein},
    description={Ein Komplex aus RNA und Proteinen. Größere Komplexe werden
    manchmal als Ribonukleoprotein-Partikel bezeichnet. (Lewin)},
    type=defi
}

\newglossaryentry{releasefaktor}{
    name={Release-Faktor},
    description={Protein, das während der Translation ein Stoppcodon erkennt und die Freisetzung der neu synthetisierten Polypeptidkette aus dem Ribosom auslöst.},
    type=defi
}

\newglossaryentry{re}{
    name={Response Element (RE)},
    plural={Response Elements (RE)},
    description={Spezifische regulatorische DNA-Sequenz, an die bestimmte Transkriptionsfaktoren als Antwort auf zelluläre Signale binden und so die Genexpression regulieren. Beispiele sind das Estrogen Response Element (ERE) oder das Glucocorticoid Response Element (GRE).},
    type=defi
}

\newglossaryentry{replikation}{
    name={Replikation},
    description={Prozess der Verdopplung der DNA vor der Zellteilung, bei dem eine identische Kopie des Genoms durch DNA-abhängige DNA-Polymerasen synthetisiert wird.},
    type=defi
}

\newglossaryentry{rev-transkription}{
    name={reverse Transkription},
    description={Synthese von DNA aus einer RNA-Matrize durch das Enzym Reverse Transkriptase, typisch für Retroviren und bestimmte mobile genetische Elemente.},
    type=defi
}


\newglossaryentry{rekombinationsfrequenz}{
    name={Rekombinationsfrequenz},
    description={Anteil rekombinanter Nachkommen an der Gesamtzahl der Nachkommen einer Kreuzung. Maß für den genetischen Abstand zweier Loci.},
    type=defi
}

\newglossaryentry{reduktionsteilung}{
    name={Reduktionsteilung},
    description={Zellteilung, bei der vor der Teilung zu einer Verrringerung (Halbierung) des Chromosomensatzes kommt},
    type=defi
}

\newglossaryentry{rapd}{
    name={Randomly Amplified Polymorphic DNA},
    description={Eine spezielle Form der Polymerase-Kettenreaktion. Es werden kurze Primer mit einer Länge von 8 bis 12 Nukleotiden verwendet, die zufällig erzeugt wurden. Aus genomischer DNA vervielfältigen sich dann nur DNA-Sequenzbereiche, die von der Sequenz des Primers eingeschlossen werden. Mit geeigneten Primern ergeben sich in der Elektrophorese individuelle Bandenmuster, die es erlauben, die Erbsubstanz unterschiedlicher Lebewesen zu vergleichen, ohne diese im Detail zu kennen. Es sind also keine aufwendigen Sequenzierungen erforderlich.},
    type=defi
}
%  \cite{docflexRAPD, dewiki:RAPD}


\newglossaryentry{punnett}{
    name={Punnett-Schema},
    description={Tabellarisches Schema zur systematischen Darstellung aller möglichen Gameten-Kombinationen einer Kreuzung und der daraus resultierenden Genotyp- und Phänotypverhältnisse.},
    type=defi
}


\newglossaryentry{purin}{
    name={Purin},
    description={Stickstoffhaltige Base mit bicyclischem Ringsystem. Zu den Purinen zählen Adenin und Guanin.},
    type=defi
}



\newglossaryentry{pyrimidin}{
    name={Pyrimidin},
    description={Stickstoffhaltige Base mit monocyclischem Ringsystem. Zu den Pyrimidinen zählen Cytosin, Thymin und Uracil.},
    type=defi
}

\newglossaryentry{pseudouridine}{
    name={Pseudouridin},
    description={Häufigste RNA-Modifikation ($\Psi$); C-glykosidische Umlagerung von Uridin, bei der C5 statt N1 die
    glykosidische Bindung trägt. Stabilisiert Sekundärstrukturen durch verbesserte Basenstapelung.},
    type=defi
}

\newglossaryentry{Proteom}{
    name={Proteom},
    description={Gesamtheit aller Proteine in einer Zelle, Gewebe, Organismus unter geg. Bedingungen (abh. von mRNA und Proteinmenge, -stabilität, -modifikationen)},
    type=defi
}

\newglossaryentry{polygene_vererbung}{
    name={Polygene Vererbung},
    description={siehe Polygenie},
    type=defi
}

\newglossaryentry{polygenie}{
    name={Polygenie},
    description={Wenn ein einzelnes Merkmal von mehreren Genen beeinflusst wird. Bspw. durch Epistasie oder
    Kopplung},
    type=defi
}

% TODO additive polygenie

\newglossaryentry{pleiotropie}{
    name={Pleiotropie},
    description={Vielfache Wirkung eines bestimmten Gens},
    type=defi
}

\newglossaryentry{phmutagenese}{
    name={physikalische Mutagenese},
    description={Mutationsauslösung durch energiereiche Strahlung: ionisierende Strahlung (Röntgen, $\gamma$) verursacht Strangbrüche und Basenmodifikationen; UV-Licht induziert Pyrimidin-Dimere.},
type=defi
}

\newglossaryentry{phänotyp}{
    name={Phänotyp},
    description={Erscheinungsbild
    eines Organismus},
    type=defi
}

\newglossaryentry{photoreparatur}{
    name={Photoreparatur},
    description={Lichtabhängiger Reparaturmechanismus, bei dem das Enzym Photolyase UV-induzierte Pyrimidindimere unter Einwirkung von sichtbarem Licht direkt spaltet.},
    type=defi
}

\newglossaryentry{oric}{
    name={oriC},
    description={Replikationsursprung des bakteriellen Chromosoms (z. B. bei E. coli), Startpunkt der bidirektionalen DNA-Replikation.},
    type=defi
}


\newglossaryentry{nondisjunction}{
    name={Nondisjunction},
    description={Das fehlende Auseinanderweichen von zwei homologen Chromosomen bei der Meiose I oder das Nichttrennen
    von Schwesterchromatiden durch eine Störung der Anaphase während der Mitose, Meiose I oder der Meiose II.},
    type=defi
}
% wikipedia

\newglossaryentry{outofafrica}{
    name={Out of Africa - Hypothese},
    description={Hypothese, dass die Gattung Homo ihren Ursprung in Afrika hatte („Wiege der Menschheit“) und dass sich
    deren Angehörige von dort über die ganze Welt verbreiteten},
    type=defi
}

\newglossaryentry{multiple_allele}{
    name={Multiple Allele},
    description={verschiedene Mutationsstufen ein und desselben Gens},
    type=defi
}

\newglossaryentry{mrna_prozessierung}{
    name={mRNA-Prozessierung},
    plural={mRNA-Prozessierungen},
    description={Gesamtheit der posttranskriptionellen Modifikationen eukaryotischer prä-mRNA, einschließlich 5'-Capping, Spleißen und Polyadenylierung, die zur reifen mRNA führen.},
    type=defi
}

\newglossaryentry{mikrotubuli}{
    name={Mikrotubuli},
    description={Röhrenförmige Elemente des Zytoskeletts bei Eukaryonten.
    Durchmesser: 24 nm},
    type=defi
}
\newglossaryentry{merkmalsform}{
    name={Merkmalsform},
    description={Konkrete Ausprägung eines Merkmals, entspricht dem Phänotyp, der durch ein bestimmtes Allel oder eine Allelkombination bedingt wird.},
    type=defi
}

\newglossaryentry{lncrna}{
    name={lncRNA},
    description={Long non-coding RNA; nicht-kodierende RNA
    $>$\,200\,nt ohne signifikantes offenes Leseraster. Reguliert
    Genexpression auf transkriptioneller (z.\,B.\ Chromatinremodellierung,
    Rekrutierung von PRC2), post-transkriptioneller und
    translationaler Ebene},
    type=defi
}

\newglossaryentry{kolinearität}{
    name={Kolinearität},
    description={Übereinstimmung der linearen Reihenfolge von Genen auf der DNA mit der Reihenfolge der genetischen Kopplungskarte beziehungsweise mit der Aminosäurereihenfolge im codierten Protein.},
    type=defi
}

\newglossaryentry{konRep}{
    name={konservative Replikation},
    description={Widerlegtes Replikationsmodell, nach dem der gesamte elterliche Doppelstrang erhalten bleibt und eine komplett neue Tochterdoppelhelix entsteht.},
    type=defi
}
\newglossaryentry{kopplungsgruppe}{
    name={Kopplungsgruppe},
    description={Gruppe von Genen, die auf demselben Chromosom liegen und daher tendenziell gemeinsam vererbt werden, da sie seltener durch Rekombination getrennt werden.},
    type=defi
}



\newglossaryentry{kinetochor}{
    name={Kinetochor},
    description={Mehrschichtiger Proteinkomplex am Centromer
    eines Chromosoms, der sich in der späten Prophase von
    Kernteilungen bildet},
    type=defi
}

\newglossaryentry{knockdown}{
    name={Knockdown-Analyse},
    description={Methode der reversen Genetik zur funktionellen Analyse
    eines Gens durch gezielte Reduktion (nicht vollständige Elimination)
    seiner mRNA-Menge.},
    type=defi
}

\newglossaryentry{kleeblatt}{
    name={Kleeblatt-Struktur},
    description={Zweidimensionale Sekundärstruktur der tRNA mit vier Doppelstrang-Helices (Akzeptor-, D-, Anticodon-
    und T$\Psi$C-Arm) und drei Schleifen; Grundlage der L-förmigen 3D-Tertiärstruktur.},
type=defi
}

\newglossaryentry{interkalation}{
    name={Interkalation},
    description={ Einlagerung planarer Moleküle zwischen benachbarte Basenpaare der DNA-Doppelhelix. Verzerrt die Helix und kann zu Leserasterverschiebungen bei der Replikation führen.},
    type=defi
}

\newglossaryentry{induzierte_mutation}{
    name={induzierte Mutation},
    description={Mutation, die durch äußere Einflüsse wie Chemikalien, Strahlung oder andere Mutagene verursacht wird.},
    type=defi
}


\newglossaryentry{Interaktom}{
    name={Interaktom},
    description={Gesamtheit aller Protein-Protein-Interaktionen},
    type=defi
}

\newglossaryentry{händigkeit}{
    name={Händigkeit},
    description={Chiralität eines Moleküls, beschreibt ob es in einer links- oder rechtsdrehenden Form vorliegt,
    relevant z. B. bei Zuckern und Aminosäuren.},
    type=defi
}


\newglossaryentry{hardy_weinberg}{
    name={Hardy-Weinberg Gesetz},
    description={Populationsgenetisches Modell, das beschreibt, unter welchen Idealbedingungen (keine Selektion, Mutation, Migration, zufällige Paarung) Allel- und Genotypfrequenzen in einer Population über Generationen konstant bleiben.},
    type=defi
}


\newglossaryentry{genkonversion}{
    name={Genkonversion},
    description={nicht-reziproke Form der Rekombination innerhalb eines Gens, bei in der ein Allel in das homologe Allel umgewandelt wird.},
    type=defi
}

\newglossaryentry{haplont}{
    name={Haplont},
    description={Organismus, dessen vorherrschende Lebensform haploid ist. Nur die Zygote ist diploid und durchläuft sofort eine Meiose.},
    type=defi
}


\newglossaryentry{genetik}{
    name={Gentik},
    description={ Die Wissenschaft, die sich mit den Gesetzen der Vererbung
    und der Variabilität der Organismen beschäftigt.},
    type=defi
}
\newglossaryentry{genotyp}{
    name={Genotyp},
    description={Gesamtheit
    der Erbanlagen (der in DNA kodierten genetischen Information) eines Organismus},
    type=defi
}

\newglossaryentry{flavonoide}{
    name={Flavonoide},
    description={  Eine Art von Farbstoffen.},
    type=defi
}


\newglossaryentry{ex-gvh}{
    name={extrinsische Genvorhersage},
    description={Methode der Genvorhersage, bei der externe biologische Informationen wie cDNA-, RNA-Seq-, EST- oder Proteinsequenzen verwendet werden, um Gene in einem Genom zu identifizieren und ihre Struktur zu bestimmen},
    type=defi
}

\newglossaryentry{excisionsreparatur}{
    name={Excisionsreparatur},
    description={DNA-Reparaturmechanismus, bei dem ein geschädigter DNA-Abschnitt herausgeschnitten und anhand des komplementären Strangs neu synthetisiert wird (Basen- oder Nukleotidexcisionsreparatur).},
    type=defi
}

\newglossaryentry{epistasie}{
    name={Epistasie},
    description={Wechselwirkung zwischen Genen an unterschiedlichen Loci, bei der die Ausprägung eines Gens durch ein oder mehrere andere Gene maskiert oder verändert wird.},
    type=defi
}

\newglossaryentry{basale_mutationsrate}{
    name={basale Mutationsrate},
    description={Häufigkeit, mit der spontane Mutationen pro Gen oder Genom und Generation ohne zusätzliche mutagene Einwirkung auftreten.},
    type=defi
}


\newglossaryentry{diplohaplont}{
    name={Diplohaplont},
    description={Bei diesen Organismen können Mitosen sowohl in der diploiden als auch in der haploiden Phase stattfinden: es entstehen zwei unterschiedlich aussehende Organismen, ein haploider und ein diploider Organismus},
    type=defi
}

\newglossaryentry{diplont}{
    name={Diplont},
    description={Organismus, dessen vorherrschende Lebensform diploid ist. Die haploide Phase beschränkt sich auf die Gameten.},
    type=defi
}


\newglossaryentry{effektorkomplex}{
    name={Effektorkomplex},
    description={Multimerer Proteinkomplex, der eine regulatorische
    RNA als Guide nutzt, um spezifische Zielmoleküle zu erkennen und
    funktionell zu beeinflussen. Im RNAi-Kontext: RISC mit AGO als
    Kernkomponente; im CRISPR-Kontext: Cas9/sgRNA-Komplex.
    Verbindet Sequenzspezifität (durch die RNA) mit katalytischer
    Aktivität (durch das Protein).},
    type=defi
}


\newglossaryentry{dsrna}{
    name={dsRNA},
    description={Doppelsträngige RNA (double-stranded RNA); entsteht
    z.\,B.\ bei Virusinfektionen, durch Transkription invertierter Repeats
    oder synthetisch. Wird von Dicer in siRNAs prozessiert und löst
    RNA-Interferenz aus. Wirkt als molekulares Gefahrensignal und
    aktiviert in Säugern zusätzlich angeborene Immunantworten},
    type=defi
}


\newglossaryentry{dnafingerprint}{
    name={DNA-Fingerabdruck},
    description={Individuenspezifisches Muster, das durch Analyse hochvariabler, repetitiver DNA-Abschnitte (z. B. VNTRs, STRs) entsteht und zur Identifizierung von Personen genutzt wird.},
    type=defi
}

\newglossaryentry{dbd}{
    name={DNA binding domain},
    description={Proteindomäne eines Transkriptionsfaktors, die sequenzspezifisch an DNA bindet und dadurch die Genregulation vermittelt.},
    type=defi
}


\newglossaryentry{disRep}{
    name={diskontinuierliche Replikation},
    description={Replikationsmodus, bei dem mindestens ein Strang, der Folgestrang, diskontinuierlich in kurzen Okazaki-Fragmenten synthetisiert wird.},
    type=defi
}
\newglossaryentry{diploid}{
    name={diploid},
    description={Auftreten von 2 homologen (mütterlichen und väterlichen) Chromosomensätzen (2 n) in den
    Kernen somatischer (2) Zellen},
    type=defi
}

\newglossaryentry{haploid}{
    name={haploid},
    description={Auftreten von nur einem Chromosomensatz (1 n) in den Zellkernen},
    type=defi
}



\newglossaryentry{dicer}{
    name={Dicer},
    description={Evolutionär konservierte RNase-III-Endonuklease, die
    lange doppelsträngige RNA (dsRNA) oder pre-miRNA in kurze
    21--23\,nt-Fragmente mit charakteristischen 2-nt-3'-Überhängen
    prozessiert (siRNA bzw.\ miRNA). Das Produkt wird anschließend
    in den RISC-Komplex geladen.},
    type=defi
}


\newglossaryentry{aneuploidie}{
    name={Aneuploidie},
    description={bweichung
    von der normalen Chromosomenzahl ganzer Chromosomensätze. (Allgemeine Genetik)},
    type=defi
}

\newglossaryentry{chmutagenese}{
    name={chemische Mutagenese},
    description={Mutationsauslösung durch chemische Agenzien},
    type=defi
}

\newglossaryentry{crossover}{
    name={Crossover},
    description={Das Ereignis, welches
    zum Austausch genetischen
    Materials zwischen Chromosomen
    führt (und dessen Resultat man
    genetisch verfolgen kann)},
    type=defi
}

\newglossaryentry{chromatid}{
    name={Chromatid},
    description={Einer der beiden Chromosomenstränge nach der Replikation des
    Chromosoms in der S-Phase des Zellzyklus},
    type=defi
}


\newglossaryentry{cohesin}{
    name={Cohesin},
    description={Proteinkomplexe, an deren Aufbau u.a. die sog. SMC Proteine
    beteiligt sind. Haben nicht nur in der Mitose eine wichtige Funktion, sie sind darüber hinaus auch an der homologen
    Rekombination in der Meiose und an der Ausbildung von Chromatinschleifen im Interphase‐Zellkern beteiligt},
    type=defi
}
% TODO besser?

\newglossaryentry{codonanticodon}{
    name={Codon-Anticodon Wechselwirkung},
    description={Antiparallele Basenpaarung zwischen dem mRNA-Codon (5'$\to$3') und dem tRNA-Anticodon (3'$\to$5')
    an der A-Stelle des Ribosoms; bestimmt die Aminosäure-Zuordnung gemäß dem genetischen Code.},
type=defi
}


\newglossaryentry{ciselement}{
    name={cis-Element},
    description={regulatorische Elemente im Promotor, ``in'' der DNA},
    type=defi
}
\newglossaryentry{codominant}{
    name={co-dominant},
    description={Vererbungsmuster, bei dem beide Allele eines heterozygoten Genotyps gleichzeitig und vollständig im Phänotyp exprimiert werden, ohne dass eines das andere maskiert.},
    type=defi
}

\newglossaryentry{chromatin}{
    name={Chromatin},
    plural={Chromatine},
    description={sichtbares ``Fadengerüst'' im
    Zellkern; Komplex aus DNA und Proteinen (vor allem Histonen), der das genetische Material in eukaryotischen Zellen
    organisiert und dessen Zugänglichkeit für Transkription und andere Prozesse reguliert.},
    type=defi
}

\newglossaryentry{chromatin_remodellierung}{
    name={Chromatin-Remodellierung},
    description={Die dynamische Anpassung der Struktur des Erbguts bei Lebewesen mit Zellkern. Es wird beispielsweise die Zugänglichkeit der genomischen DNA variiert, wodurch die Genexpression kontrolliert werden kann. Dies wird durch kovalente Histonmodifikationen mithilfe spezifischer Enzyme (z. B. Histonacetyltransferasen [HATs], Deacetylasen, Methyltransferasen und Kinasen) und durch ATP-abhängige Chromatin-Remodelling-Komplexe, die Nukleosomen entweder verschieben, ausstoßen oder restrukturieren können, gewährleistet.},
    type=defi
}
% wikipedia

\newglossaryentry{chromatinschleife}{
    name={Chromatinschleife},
    description={Schleifenförmige Struktur, in der Chromatinfasern an ein Proteingerüst (Scaffold) verankert werden. Ermöglicht die Kondensation des Chromatins auf höhere Verpackungsebenen.},
    type=defi
}



\newglossaryentry{cm}{
    name={Centimorgan},
    description={Genetische Karteneinheit; entspricht 1\,\% Rekombinationsfrequenz zwischen zwei Loci. Näherungsweise: 1\,
    cM $\approx$ 1\,Mb beim Menschen.},
type=defi
}

\newglossaryentry{anticodon}{
    name={Anticodon},
    description={Drei Nucleotide an Position 34--36 der tRNA-Sequenz, die durch komplementäre Basenpaarung mit dem mRNA-Codon
    die korrekte Aminosäure in die Peptidkette einführen.},
    type=defi
}

\newglossaryentry{asexFort}{
    name={asexuelle Fortpflanzung (vegetative, ungeschlechtliche)},
    description={ Erzeugung von Nachkommen, die identisch zu den Eltern sind.
    Die genetisch identischen Nachkommen stellen einen Klon dar.
    (Ausnahme: Mutationen)},
    type=defi
}


\newglossaryentry{biRep}{
    name={bidirektionale Replikation},
    description={Replikationsmodus, bei dem sich von einem Replikationsursprung aus zwei Replikationsgabeln in entgegengesetzte Richtungen bewegen.},
    type=defi
}



\newglossaryentry{alu_seq}{
    name={Alu-Element},
    description={Eine von mehreren verstreuten, miteinander verwandten Sequenzen, die jeweils
    etwa 300 bp lang sind, im menschlichen Genom (Mitglieder der
    SINE-Familie). Die einzelnen Mitglieder weisen an beiden Enden Alu-Schnittstellen
    auf. (Lewin)},
    type=defi
}

\newglossaryentry{aminoacyl_trna_synthasen}{
    name={Aminoacyl-tRNA-Synthetasen},
    description={Enzyme, die unter ATP-Verbrauch eine Aminosäure kovalent an das 3'-CCA-Ende der zugehörigen tRNA
    knüpfen (Aminoacylierung). Gewährleisten die korrekte Codon-Aminosäure-Zuordnung (zweiter genetischer Code).},
    type=defi
}